\section{Introduction}
Scientific papers mix prose, equations, tables, and figures, yet researchers often need to search private or unpublished PDFs without sending them to a hosted service. Retrieval-augmented generation (RAG) can expose source material at inference time~\citep{lewis2020rag}, but a deployed assistant must also distinguish a real citation from a supporting citation, a lexical match from entailment, and a refusal in final prose from an internal \textsc{Abstained} status.

We study these distinctions in \textsc{scholar}, a local-first PDF question-answering system. It ingests a document into a versioned evidence bundle, retrieves source-scoped pages and regions, answers through a local model, and records every citation, lexical support decision, edit, and terminal outcome. The intended use is researcher-assisted reading on a consumer workstation; this is an implemented prototype, not evidence of production deployment or user benefit.

The contributions are: (1) an operational design that separates acquisition-enabled setup from strict-local analysis and keeps evidence identity stable through the answer path; (2) trace-audited evaluations of retrieval fusion, deterministic post-generation repair, and post-hoc LLM-based answer assessment; and (3) adverse findings that qualify common reliability claims. In our retrospective corpus, BM25 plus dense retrieval offers little first-page gain over BM25, the tested modality heuristic harms ranking, and repair changes both citation alignment and pipeline abstention in undesirable directions. The LLM judge also identifies many incomplete answers. We report these as diagnostics, not as proof of correctness or user benefit.

The paper moves from prior scientific-document systems to ScholAR's trust boundaries and repair mechanism, then defines the evaluation before reporting results and the operational lessons they support. This order keeps implementation claims, measured outcomes, and deployment implications distinct.
